% ============================================================================
% Fichier  : partie5/P05_C20_equipe_rouge.tex
% Titre    : Équipe Rouge et Équipe Violette
% Auteur   : MINKA MI NGUIDJOÏ Thierry Emmanuel
% Version  : 1.0 -- novembre 2026
% ============================================================================
% Sources : MITRE ATT&CK ; cadre de l'évaluation fondée sur les menaces ;
% distinction test d'intrusion / équipe rouge. Opération MBAM (septembre
% 2026, BANQUE SANAGA) : cas fictif, chiffres cohérents. Deux séances
% (S23 et S24).
% ============================================================================

\chapter{Équipe Rouge et Équipe Violette}
\label{chap:rouge}
\chaptermark{Équipe rouge et équipe violette}

\epigraph{%
    « L'équipe rouge ne cherche pas toutes les failles~;\\
    elle cherche à atteindre un objectif sans se faire voir. »%
}{Principe du cours}

\niveauChapitre{Expert}{Toute la Partie~III et la
Partie~IV~; Chapitre~\ref{chap:social}.}

\begin{boxobjectifs}
\begin{itemize}
    \item Distinguer test d'intrusion, équipe rouge et équipe violette.
    \item Décrire le déroulement d'un exercice d'équipe rouge par objectifs.
    \item Comprendre la notion de scénario fondé sur une menace.
    \item Mesurer la détection et le temps de réaction du défenseur.
    \item Conduire une restitution conjointe qui améliore la défense.
\end{itemize}
\end{boxobjectifs}

\bigskip

% ============================================================================
\section*{Question directrice : la BANQUE SANAGA voit-elle une attaque discrète, et en combien de temps~?}
% ============================================================================

En septembre, avant le test de LiZbiZ, NYONG SÉCURITÉ mène pour la BANQUE
SANAGA l'Opération MBAM, un exercice d'équipe rouge par objectifs. La
question n'est plus « combien de failles~? » mais « l'attaquant atteint-il
l'objectif, et le défenseur le voit-il~? ». Cet exercice sert de référence
au cours.

\medskip
\textbf{Plan du chapitre.}
\begin{enumerate}
    \item Trois exercices distincts (section~\ref{sec:c20-distinction}).
    \item L'exercice par objectifs (section~\ref{sec:c20-objectifs}).
    \item Le scénario fondé sur une menace (section~\ref{sec:c20-menace}).
    \item Mesurer la détection (section~\ref{sec:c20-detection}).
    \item L'équipe violette (section~\ref{sec:c20-violette}).
    \item L'Opération MBAM (section~\ref{sec:c20-mbam}).
\end{enumerate}

\separateur

% ============================================================================
\section{Trois Exercices Distincts}
\label{sec:c20-distinction}
% ============================================================================

\begin{tableaucours}{L{3.2cm}L{4.6cm}L{4.4cm}}{%
    \tableauHeader{Exercice} &
    \tableauHeader{But} &
    \tableauHeader{Discrétion}
}
Test d'intrusion & Trouver le plus de failles exploitables & Peu importe \\
Équipe rouge & Atteindre un objectif précis comme le ferait un attaquant &
Essentielle \\
Équipe violette & Rouge et bleue travaillent ensemble pour améliorer la
détection & Transparente \\
\end{tableaucours}
\vspace{-0.8em}
\begin{center}{\small\itshape Tableau~20.1. Trois types d'exercice offensif}\end{center}

% ============================================================================
\section{L'Exercice par Objectifs}
\label{sec:c20-objectifs}
% ============================================================================

\begin{boxdef}
\termecle{Équipe rouge}~: exercice où une équipe offensive cherche à
atteindre un objectif convenu (un « drapeau »), par tous les moyens du
périmètre, en restant discrète, pour évaluer la capacité de détection et de
réaction de l'organisation.
\end{boxdef}

Les objectifs, ou \termecle{drapeaux}, sont définis à l'avance avec un cercle
restreint de responsables~: atteindre un système critique, obtenir un droit
d'administration, accéder à une donnée précise. Le reste de l'organisation
n'est pas prévenu, pour mesurer une détection réelle.

% ============================================================================
\section{Le Scénario Fondé sur une Menace}
\label{sec:c20-menace}
% ============================================================================

Un bon exercice d'équipe rouge imite un attaquant plausible pour
l'organisation~: ses techniques, sa patience, ses objectifs. On s'appuie sur
la connaissance des menaces et sur les techniques d'ATT\&CK pour construire
un scénario réaliste, plutôt que de dérouler des outils au hasard.

\begin{boxalert}
\textbf{Violation typique.} Transformer l'exercice en démonstration de
force qui perturbe la production.

\smallskip
\textbf{Règle~:} l'équipe rouge garde un cercle de confiance qui peut tout
arrêter, respecte des lignes à ne pas franchir (pas d'arrêt de service, pas
de données réelles exfiltrées), et consigne chaque action avec son horaire
pour la restitution.
\end{boxalert}

% ============================================================================
\section{Mesurer la Détection}
\label{sec:c20-detection}
% ============================================================================

\begin{tableaucours}{L{4.8cm}L{7.4cm}}{%
    \tableauHeader{Indicateur} &
    \tableauHeader{Ce qu'il mesure}
}
Temps jusqu'au premier objectif & Rapidité de l'attaquant \\
Temps jusqu'à la détection & Rapidité du défenseur à voir \\
Temps jusqu'à la réaction & Rapidité à contenir \\
Techniques détectées sur techniques jouées & Couverture de la détection \\
Techniques bloquées & Efficacité des parades \\
\end{tableaucours}
\vspace{-0.8em}
\begin{center}{\small\itshape Tableau~20.2. Indicateurs d'un exercice d'équipe rouge}\end{center}

% ============================================================================
\section{L'Équipe Violette}
\label{sec:c20-violette}
% ============================================================================

\begin{boxdef}
\termecle{Équipe violette}~: démarche où l'équipe offensive (rouge) et
l'équipe de défense (bleue) travaillent ensemble, en transparence~: la rouge
joue une technique, la bleue vérifie si elle la détecte, et l'on corrige
aussitôt la détection manquante.
\end{boxdef}

C'est souvent la suite la plus utile d'un exercice rouge~: on rejoue les
techniques non détectées, on écrit les règles manquantes, on vérifie qu'elles
fonctionnent. Le lab HE2-LAB-09 fait construire cette matrice de couverture.

% ============================================================================
\section{L'Opération MBAM}
\label{sec:c20-mbam}
% ============================================================================

\begin{boxscenario}{Opération MBAM. Huit jours à la BANQUE SANAGA}
\textbf{Constatations} (septembre 2026). Objectifs~: F1 (accès au réseau
interne), F2 (administration du domaine), F3 (accès au système de paiement).
\begin{itemize}
    \item F1 atteint au 2\textsuperscript{e} jour, après un hameçonnage ciblé de
          12~employés (2~clics).
    \item F2 atteint au 8\textsuperscript{e} jour (administration du domaine).
    \item F3 non atteint~: le système de paiement, cloisonné, a résisté.
    \item Détection par le SOC au 9\textsuperscript{e} jour~; confinement au 10\textsuperscript{e}.
    \item 30~techniques d'ATT\&CK jouées~: 18~détectées (dont 11 après coup
          dans les journaux), 7~bloquées.
\end{itemize}

\medskip
\textbf{Analyse} (confiance élevée). La banque a résisté là où elle avait
cloisonné (F3)~; elle a vu tard (9~jours). La couverture de détection est de
60~\% (18 sur 30), mais une part n'a été vue qu'a posteriori. La valeur de
l'exercice tient à la phase violette qui a suivi~: les 12~techniques non
détectées en direct ont donné 12~règles nouvelles.
\end{boxscenario}

\begin{boxmetier}{Analyste SOC / Red Team et RSSI}
\textbf{La vraie livraison, c'est la détection.} Un exercice d'équipe rouge
réussi ne se mesure pas au nombre de drapeaux pris, mais aux règles de
détection écrites ensuite. Le RSSI valorise la phase violette autant que
l'intrusion.
\end{boxmetier}

\separateur

\begin{boxresume}
\textbf{Ce qu'il faut retenir de ce chapitre~:}
\begin{itemize}
    \item \textbf{Test d'intrusion} (toutes les failles), \textbf{équipe
          rouge} (un objectif, discrètement), \textbf{équipe violette}
          (ensemble).
    \item \textbf{Les drapeaux} sont définis à l'avance avec un cercle
          restreint.
    \item \textbf{Le scénario} imite une menace plausible.
    \item \textbf{On mesure} temps jusqu'à l'objectif, jusqu'à la détection,
          couverture.
    \item \textbf{La phase violette} transforme l'exercice en règles de
          détection.
\end{itemize}
\end{boxresume}

\begin{boxinfo}
\textbf{Points de contrôle de l'élève.} Je sais~: distinguer les trois
exercices~; décrire un exercice par objectifs~; expliquer le scénario fondé
sur une menace~; citer les indicateurs~; expliquer la valeur de l'équipe
violette.
\end{boxinfo}

\bigskip

\begin{boxexo}[title={\ding{45}\quad Exercice~20.1. Distinction \niveaubase}]
En quoi un exercice d'équipe rouge diffère-t-il d'un test d'intrusion~?
\end{boxexo}

\begin{boxexo}[title={\ding{45}\quad Exercice~20.2. Drapeaux \niveaubase}]
Qu'est-ce qu'un drapeau~? Qui le définit~?
\end{boxexo}

\begin{boxexo}[title={\ding{45}\quad Exercice~20.3. Couverture \niveaubase}]
Calculez la couverture de détection de l'Opération MBAM.
\end{boxexo}

\begin{boxexo}[title={\ding{45}\quad Exercice~20.4. F3 \niveaumoyen}]
Pourquoi le drapeau F3 n'a-t-il pas été atteint~? Quelle leçon en tirer~?
\end{boxexo}

\begin{boxexo}[title={\ding{45}\quad Exercice~20.5. Détecté après coup \niveaumoyen}]
Quelle différence entre une technique détectée en direct et une technique
retrouvée après coup dans les journaux~?
\end{boxexo}

\begin{boxexo}[title={\ding{45}\quad Exercice~20.6. Cercle de confiance \niveaumoyen}]
Pourquoi un cercle restreint doit-il pouvoir arrêter l'exercice à tout
moment~?
\end{boxexo}

\begin{boxexo}[title={\ding{45}\quad Exercice~20.7. Lignes rouges \niveaumoyen}]
Citez trois lignes à ne pas franchir pendant un exercice d'équipe rouge.
\end{boxexo}

\begin{boxexo}[title={\ding{45}\quad Exercice~20.8. Violette \niveaumoyen}]
Décrivez le déroulement d'une séance d'équipe violette à partir des
techniques non détectées de MBAM.
\end{boxexo}

\begin{boxexo}[title={\ding{45}\quad Exercice~20.9. Scénario \niveauexpert}]
Construisez un scénario d'équipe rouge fondé sur une menace plausible pour
LiZbiZ, avec trois drapeaux.
\end{boxexo}

\begin{boxexo}[title={\ding{45}\quad Exercice~20.10. Tableau de bord \niveauexpert}]
Proposez le tableau de bord d'indicateurs que la BANQUE SANAGA devrait
suivre d'un exercice à l'autre.
\end{boxexo}

\bigskip

\begin{boxinfo}
\textbf{Pour aller plus loin.}
\begin{itemize}
    \item MITRE ATT\&CK~\cite{mitreattack}.
\end{itemize}
\end{boxinfo}

\begin{boxinfo}
\textbf{Transition.} L'exercice rouge n'a de valeur que par la défense qu'il
nourrit~: le Chapitre~\ref{chap:detection} prend le point de vue du
défenseur.
\end{boxinfo}
